% segment-local-id: OR03-S007
% segment-id: d90.orig.v1.tr03.seg0007
% license: CC BY-SA 4.0
% provenance: Materi asesmen asli berbahasa Indonesia, disusun secara independen.

\section{Set soal V: operator monoton dan metode pemisahan}
\label{orig03:ps05}

% stable-id: d90.orig.v1.tr03.ps05.exercise.0001
\begin{exercise}[Keberadaan, ketunggalan, dan titik tetap maju--mundur]
\label{orig03:ps05:exercise:0001}
Pada ruang Euklides berdimensi hingga, misalkan $f$ proper, semikontinu-bawah,
dan $\mu$-konveks kuat dengan $\mu>0$. Misalkan $g$ konveks dan
terdiferensialkan pada seluruh ruang, dengan $\nabla g$ bersifat
$\beta$-kokersif. Buktikan bahwa $f+g$ mempunyai tepat satu peminimum $x^\star$
dan bahwa, untuk setiap $\gamma>0$,
\[
 x^\star=J_{\gamma\partial f}
 \bigl(x^\star-\gamma\nabla g(x^\star)\bigr).
\]
Jelaskan peran aturan penjumlahan subdiferensial dalam argumen Anda.
\end{exercise}

% stable-id: d90.orig.v1.tr03.ps05.hint.0001
\paragraph{Petunjuk bertahap.}
\label{orig03:ps05:hint:0001}
\textbf{Tahap 1.} Gunakan batas bawah kuadratik dari konveksitas kuat untuk
mendapatkan kekompakan himpunan sublevel.
\textbf{Tahap 2.} Ubah inklusi
$0\in\partial f(x^\star)+\nabla g(x^\star)$ menjadi definisi resolven.

% stable-id: d90.orig.v1.tr03.ps05.answer.0001
\paragraph{Jawaban singkat.}
\label{orig03:ps05:answer:0001}
$f+g$ proper, semikontinu-bawah, koersif, dan $\mu$-konveks kuat; jadi ia
mempunyai peminimum tunggal. Kontinuitas $g$ memberi
$\partial(f+g)=\partial f+\nabla g$, dan syarat Fermat ekuivalen dengan
\[
 x^\star-\gamma\nabla g(x^\star)
 \in x^\star+\gamma\partial f(x^\star),
\]
yakni persamaan titik tetap yang diminta.

% stable-id: d90.orig.v1.tr03.ps05.solution.0001
\paragraph{Solusi lengkap.}
\label{orig03:ps05:solution:0001}
Konveksitas kuat memberikan batas bawah kuadratik bagi $f$ di sekitar suatu
titik dalam domainnya. Fungsi konveks terdiferensialkan $g$ mempunyai batas
bawah afin dari bidang singgungnya. Karena itu suku kuadratik mendominasi
batas afin ketika norma membesar, sehingga $f+g$ koersif. Fungsi tersebut
proper dan semikontinu-bawah, maka dalam dimensi hingga ia mencapai minimum.
Jumlahnya tetap $\mu$-konveks kuat, sehingga dua peminimum yang berbeda tidak
mungkin ada.

Karena $g$ bernilai hingga dan kontinu di seluruh ruang, aturan penjumlahan
berlaku tanpa syarat kualifikasi tambahan:
\[
 \partial(f+g)(x)=\partial f(x)+\nabla g(x).
\]
Syarat Fermat untuk peminimum tunggal adalah
$-\nabla g(x^\star)\in\partial f(x^\star)$. Mengalikan dengan $\gamma$ dan
menata ulang memberi
$x^\star-\gamma\nabla g(x^\star)\in(I+\gamma\partial f)(x^\star)$.
Menerapkan invers $J_{\gamma\partial f}=(I+\gamma\partial f)^{-1}$ memberi
persamaan titik tetap.

% stable-id: d90.orig.v1.tr03.ps05.exercise.0002
\begin{exercise}[Keaveragan operator maju--mundur dan laju residu]
\label{orig03:ps05:exercise:0002}
Misalkan $A$ maksimal monoton dan $B$ bersifat $\beta$-kokersif. Untuk
$0<\gamma<2\beta$, definisikan
\[
 T=J_{\gamma A}(I-\gamma B),
 \qquad x_{k+1}=Tx_k,
\]
dan anggap $\operatorname{Fix}T\ne\varnothing$.

\begin{enumerate}
  \item Tunjukkan bahwa $I-\gamma B$ bersifat
  $\gamma/(2\beta)$-averaged dan $J_{\gamma A}$ bersifat $1/2$-averaged.
  \item Gunakan aturan komposisi untuk memperoleh
  $\alpha=2\beta/(4\beta-\gamma)$ sebagai parameter averaged bagi $T$.
  \item Untuk $z\in\operatorname{Fix}T$, turunkan ketaksamaan penurunan dan
  batas untuk residu terkecil di antara $k$ iterasi pertama.
\end{enumerate}
\end{exercise}

% stable-id: d90.orig.v1.tr03.ps05.hint.0002
\paragraph{Petunjuk bertahap.}
\label{orig03:ps05:hint:0002}
\textbf{Tahap 1.} Tulis $I-\gamma B=(1-a)I+a(I-2\beta B)$ dengan
$a=\gamma/(2\beta)$.
\textbf{Tahap 2.} Terapkan karakterisasi kuadrat operator averaged, lalu
jumlahkan ketaksamaan itu dari $0$ sampai $k-1$.

% stable-id: d90.orig.v1.tr03.ps05.answer.0002
\paragraph{Jawaban singkat.}
\label{orig03:ps05:answer:0002}
Komposisi kedua operator tersebut bersifat
$\alpha$-averaged dengan
$\alpha=2\beta/(4\beta-\gamma)$. Untuk setiap titik tetap $z$,
\[
 \lVert Tx_k-z\rVert^2
 \le\lVert x_k-z\rVert^2
 -\frac{1-\alpha}{\alpha}\lVert(I-T)x_k\rVert^2,
\]
dan karenanya
\[
 \min_{0\le j<k}\lVert(I-T)x_j\rVert
 \le\sqrt{\frac{\alpha}{(1-\alpha)k}}\,\lVert x_0-z\rVert.
\]

% stable-id: d90.orig.v1.tr03.ps05.solution.0002
\paragraph{Solusi lengkap.}
\label{orig03:ps05:solution:0002}
Kokersivitas $B$ menyiratkan bahwa $I-2\beta B$ non-ekspansif. Dengan
$a=\gamma/(2\beta)\in(0,1)$,
\[
 I-\gamma B=(1-a)I+a(I-2\beta B),
\]
sehingga operator maju bersifat $a$-averaged. Resolven operator maksimal
monoton firmly non-ekspansif, yang sama dengan $1/2$-averaged. Jika dua
operator masing-masing $a_1$- dan $a_2$-averaged, komposisinya dapat diambil
averaged dengan parameter
\[
 \frac{a_1+a_2-2a_1a_2}{1-a_1a_2}.
\]
Memasukkan $a_1=1/2$ dan $a_2=\gamma/(2\beta)$ menghasilkan
$\alpha=2\beta/(4\beta-\gamma)\in(0,1)$.

Karakterisasi kuadrat operator $\alpha$-averaged memberi, karena $Tz=z$,
\[
 \lVert Tx_k-z\rVert^2
 \le\lVert x_k-z\rVert^2
 -\frac{1-\alpha}{\alpha}\lVert(I-T)x_k\rVert^2.
\]
Penjumlahan meneleskopkan ruas norma dan menunjukkan
\[
 \sum_{j=0}^{k-1}\lVert(I-T)x_j\rVert^2
 \le\frac{\alpha}{1-\alpha}\lVert x_0-z\rVert^2.
\]
Nilai minimum tidak melebihi rata-rata; mengambil akar memberikan batas pada
jawaban singkat dan menunjukkan residu terbaik berorde $O(k^{-1/2})$.

% stable-id: d90.orig.v1.tr03.ps05.exercise.0003
\begin{exercise}[Jendela langkah untuk operator monoton tak simetris]
\label{orig03:ps05:exercise:0003}
Untuk $a>0$, pertimbangkan operator linear
\[
 B_a=\begin{pmatrix}1&-a\\a&1\end{pmatrix}
 \quad\text{pada }\mathbb R^2.
\]
Buktikan bahwa $B_a$ bersifat $1$-monoton kuat, mempunyai konstanta Lipschitz
$\sqrt{1+a^2}$, dan bersifat $1/(1+a^2)$-kokersif. Tentukan secara eksak
kapan iterasi maju $x_{k+1}=(I-\gamma B_a)x_k$ merupakan kontraksi. Gunakan
hasil itu untuk mendiagnosis klaim bahwa monotoni kuat saja selalu membolehkan
semua $0<\gamma<2$.
\end{exercise}

% stable-id: d90.orig.v1.tr03.ps05.hint.0003
\paragraph{Petunjuk bertahap.}
\label{orig03:ps05:hint:0003}
\textbf{Tahap 1.} Untuk $d\in\mathbb R^2$, hitung
$\langle B_ad,d\rangle$ dan $\lVert B_ad\rVert^2$.
\textbf{Tahap 2.} Hitung norma faktor
$(1-\gamma)I-\gamma aJ$, dengan $J$ rotasi seperempat putaran.

% stable-id: d90.orig.v1.tr03.ps05.answer.0003
\paragraph{Jawaban singkat.}
\label{orig03:ps05:answer:0003}
Untuk setiap $d$,
$\langle B_ad,d\rangle=\lVert d\rVert^2$ dan
$\lVert B_ad\rVert^2=(1+a^2)\lVert d\rVert^2$. Faktor kontraksinya adalah
\[
 \rho(\gamma)=\sqrt{(1-\gamma)^2+a^2\gamma^2},
\]
sehingga kontraksi terjadi tepat ketika
$0<\gamma<2/(1+a^2)$.

% stable-id: d90.orig.v1.tr03.ps05.solution.0003
\paragraph{Solusi lengkap.}
\label{orig03:ps05:solution:0003}
Bagian simetris $B_a$ adalah identitas. Karena itu
$\langle B_ad,d\rangle=\lVert d\rVert^2$, yang membuktikan monotoni kuat
dengan modulus satu. Perkalian langsung memberi
$B_a^\top B_a=(1+a^2)I$; jadi norma operatornya
$\sqrt{1+a^2}$ dan
\[
 \langle B_ad,d\rangle
 =\frac1{1+a^2}\lVert B_ad\rVert^2,
\]
yakni kokersivitas yang dinyatakan.

Matriks iterasi mempunyai bagian skalar $1-\gamma$ dan bagian rotasi
$-\gamma aJ$, sehingga norma setiap vektor dikalikan dengan
$\rho(\gamma)$. Syarat $\rho(\gamma)<1$ ekuivalen dengan
\[
 (1-\gamma)^2+a^2\gamma^2<1
 \quad\Longleftrightarrow\quad
 0<\gamma<\frac{2}{1+a^2}.
\]
Jadi interval $(0,2)$ tidak seragam terhadap bagian antisimetri. Monotoni kuat
tanpa kontrol Lipschitz atau kokersivitas tidak cukup untuk membenarkan klaim
langkah tersebut.

% stable-id: d90.orig.v1.tr03.ps05.exercise.0004
\begin{exercise}[Sertifikat titik tetap Douglas--Rachford]
\label{orig03:ps05:exercise:0004}
Pada $\mathbb R$, minimalkan $f(x)+g(x)$ dengan
$f(x)=|x|$ dan $g=\delta_{[1,\infty)}$. Ambil $\gamma=1$ dan konvensi
\[
 T=\tfrac12(I+R_gR_f),
 \qquad R_q=2\operatorname{prox}_q-I.
\]
Tentukan peminimum, lalu verifikasi secara langsung bahwa $z^\star=2$ adalah
titik tetap $T$ dan bahwa bayangannya
$\operatorname{prox}_f(z^\star)$ sama dengan peminimum.
\end{exercise}

% stable-id: d90.orig.v1.tr03.ps05.hint.0004
\paragraph{Petunjuk bertahap.}
\label{orig03:ps05:hint:0004}
\textbf{Tahap 1.} Gunakan penyusutan lunak untuk $\operatorname{prox}_{|\cdot|}$
dan proyeksi untuk $\operatorname{prox}_{\delta_{[1,\infty)}}$.
\textbf{Tahap 2.} Ikuti urutan $2\mapsto R_f(2)\mapsto R_gR_f(2)$ tanpa
menukar konvensi komposisi.

% stable-id: d90.orig.v1.tr03.ps05.answer.0004
\paragraph{Jawaban singkat.}
\label{orig03:ps05:answer:0004}
Peminimumnya $x^\star=1$. Selain itu,
$\operatorname{prox}_f(2)=1$, $R_f(2)=0$,
$\operatorname{prox}_g(0)=1$, dan $R_g(0)=2$. Maka $T(2)=2$ dan bayangan
$\operatorname{prox}_f(2)=x^\star$.

% stable-id: d90.orig.v1.tr03.ps05.solution.0004
\paragraph{Solusi lengkap.}
\label{orig03:ps05:solution:0004}
Kendala efektif $g$ memaksa $x\ge1$. Pada himpunan itu $|x|=x$, sehingga
nilai terkecil dicapai secara tunggal di $x^\star=1$. Proksimal $f$ dengan
parameter satu adalah penyusutan lunak; akibatnya
$\operatorname{prox}_f(2)=2-1=1$ dan $R_f(2)=2(1)-2=0$. Proksimal indikator
adalah proyeksi ke $[1,\infty)$, sehingga
$\operatorname{prox}_g(0)=1$ dan $R_g(0)=2(1)-0=2$. Oleh sebab itu
\[
 T(2)=\tfrac12(2+2)=2.
\]
Variabel titik tetap $z^\star$ tidak harus sama dengan solusi primal; pada
konvensi ini solusi diperoleh dari bayangan proksimalnya, yang memang bernilai
satu.

% stable-id: d90.orig.v1.tr03.ps05.exercise.0005
\begin{exercise}[Metode titik proksimal dengan galat yang dapat dijumlahkan]
\label{orig03:ps05:exercise:0005}
Misalkan $A$ maksimal monoton pada ruang Euklides berdimensi hingga,
$\operatorname{zer}A\ne\varnothing$, dan $\gamma>0$. Pertimbangkan
\[
 x_{k+1}=J_{\gamma A}x_k+e_k,
 \qquad \sum_{k=0}^{\infty}\lVert e_k\rVert<\infty.
\]
Buktikan bahwa $(x_k)$ konvergen ke suatu titik di $\operatorname{zer}A$.
Dalam bukti Anda, tunjukkan juga bahwa
$x_k-J_{\gamma A}x_k\to0$.
\end{exercise}

% stable-id: d90.orig.v1.tr03.ps05.hint.0005
\paragraph{Petunjuk bertahap.}
\label{orig03:ps05:hint:0005}
\textbf{Tahap 1.} Tetapkan $y_k=J_{\gamma A}x_k$ dan gunakan
firm non-ekspansivitas terhadap $z\in\operatorname{zer}A$.
\textbf{Tahap 2.} Jumlahkan galat kuadrat yang dihasilkan; lalu gunakan
ketertutupan graf operator maksimal monoton pada subsekuens konvergen.

% stable-id: d90.orig.v1.tr03.ps05.answer.0005
\paragraph{Jawaban singkat.}
\label{orig03:ps05:answer:0005}
Firm non-ekspansivitas memberi
\[
 \lVert y_k-z\rVert^2+\lVert x_k-y_k\rVert^2
 \le\lVert x_k-z\rVert^2.
\]
Perturbasi yang normanya dapat dijumlahkan membuat barisan kuasi-Fej\'{e}r,
dan penjumlahan ketaksamaan memperlihatkan
$\sum_k\lVert x_k-y_k\rVert^2<\infty$. Semua titik klaster adalah nol $A$;
dalam dimensi hingga sifat kuasi-Fej\'{e}r lalu memberi konvergensi seluruh
barisan.

% stable-id: d90.orig.v1.tr03.ps05.solution.0005
\paragraph{Solusi lengkap.}
\label{orig03:ps05:solution:0005}
Ambil $z\in\operatorname{zer}A$. Maka $J_{\gamma A}z=z$. Dengan
$y_k=J_{\gamma A}x_k$, firm non-ekspansivitas menghasilkan
\[
 \lVert y_k-z\rVert^2+\lVert x_k-y_k\rVert^2
 \le\lVert x_k-z\rVert^2.                                      \tag{1}
\]
Karena $x_{k+1}=y_k+e_k$,
$\lVert x_{k+1}-z\rVert\le\lVert x_k-z\rVert+\lVert e_k\rVert$.
Jumlah norma galat hingga, sehingga $(x_k)$ dan $(y_k)$ terbatas dan jarak
terhadap setiap nol mempunyai limit.

Mengembangkan kuadrat $\lVert y_k+e_k-z\rVert^2$ lalu memakai (1) memberi
\[
 \lVert x_{k+1}-z\rVert^2
 \le\lVert x_k-z\rVert^2-\lVert x_k-y_k\rVert^2
   +2\langle y_k-z,e_k\rangle+\lVert e_k\rVert^2.
\]
Karena $(y_k)$ terbatas dan galat dapat dijumlahkan, kedua suku galat di ruas
kanan dapat dijumlahkan. Penjumlahan ketaksamaan menunjukkan
$\sum_k\lVert x_k-y_k\rVert^2<\infty$, jadi residu itu menuju nol.

Setiap subsekuens konvergen $x_{k_j}\to\bar x$ juga memenuhi
$y_{k_j}\to\bar x$. Dari definisi resolven,
$(x_{k_j}-y_{k_j})/\gamma\in Ay_{k_j}$, dan ruas kiri menuju nol. Graf
operator maksimal monoton tertutup, sehingga $0\in A\bar x$. Jadi semua titik
klaster berada di $\operatorname{zer}A$. Barisan terbatas dalam dimensi hingga
mempunyai titik klaster, dan sifat kuasi-Fej\'{e}r beserta fakta bahwa semua
titik klasternya adalah nol memaksa seluruh barisan konvergen ke satu nol
$A$.
